\chapter{Non-abelian Poincaré duality and the Verdier dual of the bar coalgebra}
\label{chap:NAP-koszul-derivation}

\index{Poincare duality!chiral|textbf}
\index{Verdier duality!bar coalgebra and dual algebra}
\index{Theorem C!three-tier structure C0, C1, C2}

\begin{remark}[Theorem~C three-tier orientation: C0, C1, C2]
\label{rem:theorem-c-three-tier-orientation}
\index{Theorem C!C0 fiber-center, C1 Lagrangian eigenspace, C2 conditional}
Complementarity (Theorem~C) decomposes into three tiers, never
to be conflated.
\textsc{C0}: the fibre-centre identification
$Z(\cA)=H^0Z^{\mathrm{der}}_{\mathrm{ch}}(\cA)$ with
$Z(\Omegach_X(\Bbarch_X(\cA)))$ at every chirally Koszul datum,
proved on that datum by the bar-cobar counit
$\Omegach_X\Bbarch_X(\cA)\xrightarrow{\sim}\cA$
of Theorem~\ref{thm:koszul-reflection}.
\textsc{C1}: the Lagrangian-eigenspace decomposition of the
quantum class
\[
Q_g(\cA) \oplus Q_g(\cA^!)
\simeq H^*(\overline{\cM}_g,\mathcal{Z}_{\cA})
\]
at every genus $g \geq 1$ on the finite-type or completed Koszul
locus where the bar--cobar normal complex is perfect and the
Verdier pairing is dualizable. This is proved family-by-family
on the standard landscape; outside that landscape it is conditional
on perfectness, separated completion, proper support, and the chosen
orientation datum.
\textsc{C2}: the upgrade to a $(-1)$-shifted symplectic Lagrangian
intersection of $\cM_\cA$ and $\cM_{\cA^!}$ inside the global
deformation space $\cM_{\mathrm{comp}}$, conditional on the BV
package and the bar-chart Lagrangian lift; the uniform-weight
restriction is \emph{not} a hypothesis of \textsc{C2} but enters
only at the scalar-complementarity corollary
$\kappa(\cA) + \kappa(\cA^!) = K(\cA)$. The scalar conductor
$K(\cA) = \kappa(\cA) + \kappa(\cA^!)$ is the trace shadow of
\textsc{C1} together with Theorem~D, and the Climax identification
$\kappa(\cA) = -c_{\mathrm{ghost}}(\BRST(\cA))$
(Chapter~\ref{chap:climax-platonic};
Chapter~\ref{chap:universal-conductor}) factors $K(\cA)$ as the
universal ghost-charge sum across the Koszul pair.
\end{remark}

The MC element $\Theta_\cA$ lives in a convolution algebra built
from two objects only: the bar coalgebra $\barB_X(\cA)$ and the
algebra~$\cA$. The Verdier dual of the bar coalgebra produces a
third object,
\begin{equation}\label{eq:homotopy-koszul-dual-def}
\cA^!_\infty \;:=\; \mathbb{D}_{\Ran}\barB_X(\cA),
\end{equation}
a factorization algebra on $\Ran(X)$. On the finite Koszul locus,
when the transferred post-Verdier $A_\infty$-structure is formal,
this algebra has strict model
$\cA^!=(\cA^i)^\vee$, where
$\cA^i=H^\bullet(\barB_X(\cA))$ is the Koszul-dual coalgebra.
The cobar counit
$\Omegach_X\barB_X(\cA)\xrightarrow{\sim}\cA$ is bar-cobar
inversion, not Koszul duality.

Verdier duality is used on the bounded holonomic, equivalently
constructible, Verdier-dualizable subcategory; in the completed
bar-length setting it means the continuous dual on the
pro-conilpotent completion. It is not asserted on an arbitrary
unbounded category of $\mathcal{D}$-modules. On this surface
$\mathbb{D}_{\Ran} \colon
D^b_{\mathrm{hol}}(\mathcal{D}\text{-mod}(\Ran(X)))
\to D^b_{\mathrm{hol}}(\mathcal{D}\text{-mod}(\Ran(X)))$
exchanges $!$-tensor and $*$-tensor. Hence it carries
factorization coalgebras to factorization algebras:
\begin{center}
\begin{tabular}{r@{\;$\longleftrightarrow$\;}l}
bar coproduct $\Delta_{\barB}$ & chiral product $\mu$ of $\cA^!_\infty$, \\
bar counit $\varepsilon_{\barB}$ & unit of $\cA^!_\infty$, \\
bar differential $d_{\barB}$ & $A_\infty$-differential of $\cA^!_\infty$, \\
logarithmic forms on $\overline{C}_n(X)$ & distributional
 currents on $C_n(X)$.
\end{tabular}
\end{center}
No pre-existing dual partner, quadratic presentation, or
classical $R^\perp$ is part of the definition.

The firewall is:
\begin{center}
\begin{tabular}{r@{\;:\quad}l}
$\cA$ & chiral algebra, the boundary object, \\
$\barB_X(\cA)$ & factorization coalgebra, the bar chain object, \\
$\cA^i=H^\bullet(\barB_X(\cA))$ & Koszul-dual coalgebra on the
 finite or completed Koszul locus, \\
$\cA^! = (\cA^i)^\vee$ & Verdier or continuous linear dual algebra,
 after finite-type or completed dualizability, \\
$Z^{\mathrm{der}}_{\mathrm{ch}}(\cA)
 = \ChirHoch^\bullet(\cA,\cA)$ & cochain-level chiral closed sector, not a
 bar object and not topological or categorical Hochschild.
\end{tabular}
\end{center}
The genus-$g$ obstruction theory uses the post-Verdier algebra
$\cA^!_\infty$ and its strict model~$\cA^!$. The cochain-level
closed-sector object is $Z^{\mathrm{der}}_{\mathrm{ch}}(\cA)$; a
physical bulk realization requires an additional open-closed comparison
and analytic completion.

\section{The intrinsic duality datum}

\begin{problem}[Independent construction of Koszul dual]\label{prob:independent-dual}
An identification of the bar coalgebra of $\mathcal{A}_1$ with the
algebra $\mathcal{A}_2^!$ is typed only after the Verdier duality
step is included.
The left side is a factorization coalgebra. The right side is a
factorization algebra. The correct pair of statements is
\[
H^\bullet\bar{B}^{\mathrm{ch}}(\mathcal{A}_1)
\simeq \mathcal{A}_1^i
\qquad\text{and}\qquad
\mathbb{D}_{\Ran}\bar{B}^{\mathrm{ch}}(\mathcal{A}_1)
\simeq \mathcal{A}_2\simeq\mathcal{A}_1^!
\]
for a finite-type Koszul pair after the Verdier-compatible
identification, with continuous completions off it.
The completed statement requires a separated complete
bar filtration, finite-dimensional graded quotients, and continuous
duals; without these hypotheses no algebra object is defined.
The cobar equivalence
$\Omega^{\mathrm{ch}}\bar{B}^{\mathrm{ch}}(\mathcal{A}_1)
\simeq \mathcal{A}_1$
is reconstruction of the original algebra, not construction of
the dual algebra.
\end{problem}

\begin{remark}[NAP duality]\label{rem:NAP-solution}
Non-abelian Poincar\'e duality supplies the post-Verdier algebra.
Let $M$ be compact and oriented, and let $\mathcal{A}$ be an
augmented locally constant or constructible factorization algebra
whose local values are dualizable and whose global sections satisfy
the required compact-support finiteness. If its Verdier-linear
Koszul dual algebra $\mathcal{A}^!$ exists in the finite-type or
completed continuous sense, then Ayala--Francis duality gives:
\[
\int_M \mathcal{A}
\;\simeq\;
\mathbb{D}\!\left(\int_{-M} \mathcal{A}^!\right).
\]
Here $\int_M$ is factorization homology, $\mathbb{D}$ is Verdier
duality, $-M$ is $M$ with opposite orientation, and
$\mathcal{A}^!$ denotes the algebra obtained after Verdier or
continuous linear duality, not the bar coalgebra. The formula is a
theorem under these compactness, orientation, support, and
dualizability hypotheses; it is not a definition for an arbitrary
chiral algebra.

Factorization homology is computed as a colimit over Weiss covers:
\[\int_M \mathcal{A} = \operatorname{colim}_{\coprod_i D_i \hookrightarrow M} \bigotimes_i \mathcal{A}(D_i)\]
where the colimit ranges over collections of disjoint disk embeddings $\coprod_i D_i \hookrightarrow M$ (not over arbitrary configuration spaces; cf.\ the excision property for Weiss covers, Definition~\ref{def:weiss-cover}).

Verdier duality exchanges integration over $\overline{C}_n(M)$
with distributions on $C_n(M)$, logarithmic forms with delta
currents, and residues at collisions with singularity insertions.
Construction~\ref{const:A-dual-intrinsic} fixes the chiral
object-level version of this equivalence.
\end{remark}

\section{Logarithmic forms and distributional currents}
\label{sec:geometric-arena}
\index{collision geometry|textbf}

Configuration spaces carry the duality that separates the bar and
cobar sides. On $\overline{C}_k(X)$, two dual sheaf complexes
coexist; one builds the bar construction, the other the cobar.

The sheaf $\Omega^*_{\log}(\overline{C}_k(X))$ consists of
differential forms with at worst logarithmic poles along the
boundary divisors $D_I = \{z_i = z_j : i, j \in I\}$. The
propagator $\eta_{ij} = d\log(z_i - z_j)$ is the simplest
section: it has a simple pole on $D_{ij}$ and is smooth elsewhere.
The wedge products $\bigwedge_{e \in E(T)} \eta_e$, indexed by edges
of planar trees, are the integrands of the bar construction.

The notation $\mathcal{D}^*_{\mathrm{dist}}(C_k(X))$ denotes the
de~Rham representative of the Verdier-dual constructible complex:
compactly supported currents on the open stratum together with the
boundary-supported dualizing sections on $\overline{C}_k(X)$ that
arise from the extension triangles. The delta symbol
$\delta_{ij}$ is the current representative of the localization class
along the boundary divisor $D_{ij}$, not a function on $C_k(X)$.
Products such as $\prod_{e\in E(T)}\delta_e$ are names for iterated
localization correspondences, with derived intersection understood;
they are not pointwise products of singular distributions.

Verdier duality exchanges these two constructible complexes. After
compact support on the open stratum and finite-dimensional
cohomology are imposed, the de~Rham notation above realizes the
hypercohomological pairing
\begin{equation}\label{eq:log-dist-pairing-pd}
\langle \cdot, \cdot \rangle \colon
R\Gamma(\overline{C}_k(X),\Omega^\bullet_{\log}(D))
\otimes
R\Gamma_c(C_k(X),\mathcal{D}^\bullet_{\mathrm{dist}})
\longrightarrow \mathbb{C}
\end{equation}
where $j\colon C_k(X)\hookrightarrow\overline{C}_k(X)$ is the open
inclusion. A current representative $K$ acts by the regularized
evaluation $K(j^*\omega)$, equivalently by the residue sum over the
boundary strata. No pullback of a boundary current to the open
stratum and no pointwise product of singular kernels is part of the
definition. Residue extraction (bar) and singularity insertion
(cobar) are Poincar\'e--Lefschetz dual operations on the pair
$(\overline{C}_k(X),D)$.

\begin{example}[Two points]
For $k = 2$ on $X = \mathbb{P}^1$:
$\langle \eta_{12}, \delta_{12} \rangle
= \operatorname{Res}_{z_1 = z_2} \eta_{12} = 1$.
The OPE coefficient of the identity channel is the residue of the
propagator, paired with the delta current at collision.
\end{example}

\begin{example}[Three points and the Arnold relation]
For $k = 3$, the three boundary divisors $D_{12}$, $D_{23}$,
$D_{31}$ of $\overline{C}_3(\mathbb{P}^1)$ satisfy
$\eta_{12} \wedge \eta_{23} + \eta_{23} \wedge \eta_{31}
+ \eta_{31} \wedge \eta_{12} = 0$
(Theorem~\ref{thm:arnold-three}). Pairing the corresponding
constructible diagonal classes with these logarithmic forms gives
the associativity constraint: the signed sum of iterated residues
vanishes. The literal product of delta distributions is only a local
mnemonic; Remark~\ref{rem:analytic-framework} gives the derived
formulation.
\end{example}

\section{Verdier duality on configuration spaces}

\begin{setup}[Configuration space duality]\label{setup:config-verdier}
Let $X$ be a smooth curve (or more generally, an $n$-dimensional manifold). The configuration space of $k$ points is:
\[C_k(X) = \{(z_1, \ldots, z_k) \in X^k : z_i \neq z_j \text{ for } i \neq j\}\]

Its Fulton--MacPherson compactification $\overline{C}_k(X)$ is a smooth manifold with corners, with boundary divisors parametrizing collision patterns.
\end{setup}

\begin{theorem}[Verdier duality for configuration spaces; \ClaimStatusProvedElsewhere{} \cite{KS90}]\label{thm:verdier-config}
\index{Verdier duality!configuration spaces}
Let $X$ be a smooth oriented curve, let
$j\colon C_k(X)\hookrightarrow \overline{C}_k(X)$ be the open
configuration stratum of its Fulton--MacPherson compactification,
and let $D=\overline{C}_k(X)\setminus C_k(X)$ be the normal-crossing
boundary. For the logarithmic de~Rham model of
$Rj_*\mathbb C_{C_k(X)}$ and the compactly supported current model
of its Verdier dual, Verdier duality gives a canonical pairing
\[
R\Gamma(\overline{C}_k(X),\Omega^\bullet_{\log}(D))
\otimes
R\Gamma_c(C_k(X),\mathcal{D}^\bullet_{\mathrm{dist}})
\longrightarrow \mathbb{C}.
\]
On finite-dimensional cohomology this pairing is perfect. In a
de~Rham model it is represented by the action of a current on the
regularized logarithmic form,
\[
\langle \omega, K \rangle = K(j^*\omega),
\]
or, after resolving the boundary stratification, by the corresponding
sum of residue pairings on the strata.
The Verdier-dual identification is an isomorphism in
$D^b_c(\overline{C}_k(X))$ with the orientation shift
$d=\dim_{\mathbb{R}}\overline{C}_k(X)$; it is not a
degree-by-degree isomorphism of raw logarithmic forms with raw
distributions.
\end{theorem}

\begin{proof}
We apply Verdier duality for constructible
sheaves~\cite{KS90} to the specific geometry of
configuration spaces.

\emph{Step 1.}
For any constructible complex $\mathcal{F}$ on $\overline{C}_k(X)$:
\[\mathbb{D}(\mathcal{F}) = \mathcal{RHom}_{\overline{C}_k(X)}(\mathcal{F}, \omega_{\overline{C}_k(X)}[d])\]

\emph{Step 2.}
Since $\overline{C}_k(X)$ is smooth of complex dimension $k$ (for $X$ a smooth curve), the constant sheaf is self-dual up to shift:
\[\mathbb{D}(\mathbb{C}_{\overline{C}_k(X)}) \cong \mathbb{C}_{\overline{C}_k(X)}[2k]\]
For the $j_!$-extension from the open stratum $C_k(X) \hookrightarrow \overline{C}_k(X)$, Verdier duality exchanges $j_!$ and $Rj_*$:
\[\mathbb{D}(j_! \mathbb{C}_{C_k(X)}) \cong Rj_* \mathbb{C}_{C_k(X)}[2k]\]

\emph{Step 3.}
Taking hypercohomology with respect to the de Rham complex:
\begin{align*}
\mathbb{H}^*(\overline{C}_k(X), \Omega^*_{\log}) &= H^*_{\text{dR}}(\overline{C}_k(X), \log D)\\
\mathbb{H}^*_{C_k(X)}(\Omega^*_{\text{dist}}) &= H^*_{c,\text{dR}}(C_k(X))
\end{align*}

The Verdier duality pairing descends to the de Rham pairing.

\emph{Step 4.}
The explicit pairing is computed by residues:
\[\langle \omega_{\text{log}}, K_{\text{dist}} \rangle = \sum_{\text{strata } S} \int_S \text{Res}_S(\omega_{\text{log}}) \wedge K_{\text{dist}}|_S\]

Verdier duality for constructible complexes on smooth varieties~\cite{KS90}
gives perfectness after passing to finite-dimensional cohomology.
\end{proof}

\begin{remark}[Algebraic vs.\ analytic Verdier duality]\label{rem:verdier-bridge}
\index{Verdier duality!algebraic vs.\ analytic}
Theorem~\ref{thm:verdier-config} uses ``Verdier duality'' in two senses that
must be carefully distinguished:
\begin{enumerate}
\item \emph{Algebraic ($\mathcal{D}$-module) Verdier duality.}
The functor $\mathbb{D}_{\mathrm{alg}} = \mathcal{RH}om(-,
\omega_{\overline{C}_k}[\dim])$ acts on holonomic
$\mathcal{D}$-modules on $\overline{C}_k(X)$. It exchanges
$j_!$ and $j_*$ extensions (Lemma~\ref{lem:verdier-extension-exchange}).
This is the sense used in Definition~\ref{def:geom-cobar-intrinsic} to
construct the cobar complex intrinsically.
\item \emph{Analytic (Borel--Moore) Verdier duality.}
On the real manifold underlying $C_k(X)$, Poincar\'e--Lefschetz
duality identifies $H^*(\overline{C}_k)$ with $H^{d-*}_c(C_k)$,
realized concretely by the integration pairing of logarithmic forms
against distributional currents.
\end{enumerate}
The bridge between (1) and (2) is the Riemann--Hilbert correspondence:
applying the de~Rham functor to holonomic $\mathcal{D}$-modules produces
constructible sheaves, and the algebraic duality $\mathbb{D}_{\mathrm{alg}}$
becomes the topological Verdier duality on constructible complexes
\cite{KS90}. See Theorem~\ref{thm:cobar-distributional-model}
for the analytic realization as a consequence of the algebraic statement.
\end{remark}

\subsection{The dual operations}

\begin{theorem}[Dual differentials; \ClaimStatusProvedHere]\label{thm:dual-differentials}
In the finite-type constructible model of
Theorem~\ref{thm:verdier-config}, the following operations are
Verdier-adjoint:

\emph{Residue vs.\ boundary localization.}
Let $i_{ij}\colon D_{ij}\hookrightarrow\overline{C}_k(X)$ be the
boundary inclusion and
$q_{ij}\colon D_{ij}\to\overline{C}_{k-1}(X)$ the collision
collapse map. The cobar operation denoted by $\delta_{ij}$ is the
localization correspondence $i_{ij,*}q_{ij}^{!}$.
\begin{align*}
\text{Bar: } &\text{Res}_{D_{ij}}: \Omega^*_{\log}(\overline{C}_k) \to \Omega^*_{\log}(\overline{C}_{k-1})\\
\text{Cobar: } &\delta_{ij}: \mathcal{D}^*_{\text{dist}}(C_{k-1}) \to \mathcal{D}^*_{\text{dist}}(C_k)
\end{align*}

Explicitly:
\[\langle \text{Res}_{D_{ij}}(\omega), K \rangle = \langle \omega, \delta_{ij}(K) \rangle\]

\emph{Collapsing vs.\ splitting.}
\begin{align*}
\text{Bar: } &\text{Collapse at } D: C_k \dashrightarrow C_{k-1}\\
\text{Cobar: } &\text{Split along diagonal: } C_{k-1} \to C_k
\end{align*}

\emph{Composition product vs.\ coproduct.}
\begin{align*}
\text{Bar: } &\circ: \Omega^*_{\log}(\overline{C}_k) \times \Omega^*_{\log}(\overline{C}_l) \to \Omega^*_{\log}(\overline{C}_{k+l-1})\\
\text{Cobar: } &\Delta: \mathcal{D}^*_{\text{dist}}(C_{k+l-1}) \to \mathcal{D}^*_{\text{dist}}(C_k) \otimes \mathcal{D}^*_{\text{dist}}(C_l)
\end{align*}
\end{theorem}

\begin{proof}[Geometric proof via stratifications]
\emph{Observation \textup{(Serre)}.} Boundary divisors $D \subset \overline{C}_k(X)$ biject with collision patterns (combinatorial data) and with diagonal subspaces in $C_k(X)$ (geometric data).

\emph{Residue operation.}
For a form $\omega$ with logarithmic pole along $D$:
\[\omega = f \cdot \eta_{ij} \wedge \omega' + \text{regular}\]
where $\eta_{ij} = d\log(z_i - z_j)$.

The residue is:
\[\text{Res}_D(\omega) = f|_D \cdot \omega'|_D\]

\emph{Boundary-localization operation.}
For a current representative $K$ on $C_{k-1}(X)$, the Verdier-dual
operation is
\[
\delta_{ij}(K):=i_{ij,*}q_{ij}^{!}K.
\]
In local analytic notation this class is often written as
$K(z_1,\ldots,\widehat z_i,\ldots,z_k)\delta(z_i-z_j)$, but the
display is only the current representative of the localization
morphism.

\emph{The duality.}
The following calculation is the local de~Rham model of the
Verdier adjunction; the product with $\delta(z_i-z_j)$ is shorthand
for the localization morphism supported on the diagonal:
\begin{align*}
\langle \omega, \delta_{ij}(K) \rangle &= \int_{\overline{C}_k} (f \cdot \eta_{ij} \wedge \omega') \wedge \delta(z_i - z_j) \cdot K\\
&= \int_{D} f|_D \cdot \omega'|_D \wedge K|_D \quad \text{(localization to the boundary)}\\
&= \int_{\overline{C}_{k-1}} \text{Res}_D(f \cdot \eta_{ij} \wedge \omega') \wedge K\\
&= \langle \text{Res}_D(\omega), K \rangle
\end{align*}
\end{proof}

\begin{example}[Heisenberg specialization]\label{ex:dual-diff-heisenberg}
For the Heisenberg algebra $\cH_k$ with generator $J(z)$,
Theorem~\ref{thm:dual-differentials} reduces to:
the residue $\operatorname{Res}_{z_1=z_2} J(z_1) J(z_2)\, d\log(z_1-z_2) = k$
on the bar side is dual to the boundary-localization insertion
$\delta(z_1-z_2)$ on the cobar side, recovering the
level-$k$ pairing that identifies $\cH_k^! = \mathrm{Sym}^{\mathrm{ch}}(V^*)$
in \S\ref{sec:frame-koszul-dual}.
\end{example}

\begin{remark}[Analytic framework]\label{rem:analytic-framework}
The computation above is formal: the expression
$\eta_{ij} \wedge \delta(z_i - z_j)$ is not well-defined as a
pointwise product of distributions, since
$\eta_{ij}=d\log(z_i-z_j)$ has a singularity precisely at the
support of $\delta(z_i-z_j)$. The correct framework is Verdier
duality as a functor on constructible sheaves
(see \cite[Expos\'e XVIII]{SGA4}). The pairing is defined in
$D^b_c(\overline{C}_k(X))$ by proper push-forward and the
localization triangle for $(C_k(X),D_{ij})$, not by multiplying a
delta current by a singular logarithmic form.
\end{remark}

\section{From Verdier duality to cooperad structure}

\begin{construction}[Post-Verdier homotopy dual algebra]\label{const:A-dual-intrinsic}
Let $\mathcal{A}$ be an augmented chiral algebra on a smooth
projective oriented curve $X$, valued in regular holonomic
$\mathcal{D}$-modules. Assume each Fulton--MacPherson component is
proper over the relevant Ran stratum and Verdier-dualizable; if the
bar-length filtration is infinite, work in the separated
pro-conilpotent completion and use continuous duals. The geometric
bar construction is the factorization coalgebra
\[
\barB_X^{\mathrm{ch}}(\mathcal{A})
=
\bigoplus_{n\geq 1}
\pi_{n,*}\!\left(
j_*j^*(\mathcal{A}^{\boxtimes n})
\otimes \Omega^\bullet_{\log}(\overline{C}_n(X))
\right),
\]
with coproduct induced by splitting configurations and differential
given by internal differential, de~Rham differential, and Poincar\'e
residue along collision divisors.

Define the post-Verdier homotopy dual algebra by
\[
\mathcal{A}^!_\infty
:=
\mathbb{D}_{\Ran}\barB_X^{\mathrm{ch}}(\mathcal{A}).
\]
Its $n$-ary component is represented by distributional sections on
$C_n(X)$:
\[
(\mathcal{A}^!_\infty)_n
\simeq
\mathbb{D}_{\overline{C}_n(X)}
\!\left(
j_*j^*(\mathcal{A}^{\boxtimes n})
\otimes \Omega^\bullet_{\log}(\overline{C}_n(X))
\right).
\]
The bar coproduct dualizes to the product of
$\mathcal{A}^!_\infty$, the bar counit dualizes to the unit, and
the bar differential dualizes to the $A_\infty$ differential.

On the finite Koszul locus, where the bar cohomology is finite in
each graded piece,
\[
\mathcal{A}^i := H^\bullet\barB_X^{\mathrm{ch}}(\mathcal{A})
\]
is the Koszul-dual coalgebra and
\[
\mathcal{A}^! := (\mathcal{A}^i)^\vee
\]
is the strict Koszul-dual algebra. The comparison
$\mathcal{A}^!_\infty\simeq\mathcal{A}^!$ is formality of the
post-Verdier algebra, not a definition of $\mathcal{A}^!$ as a
coalgebra.
\end{construction}

\begin{remark}\label{rem:intrinsic-nature}
Construction~\ref{const:A-dual-intrinsic} uses no pre-existing dual
partner and no orthogonal-relation presentation. The chain coalgebra
is $\barB_X(\cA)$. The cohomological coalgebra is
$\cA^i=H^\bullet\barB_X(\cA)$. The algebra is
$\cA^!_\infty=\mathbb{D}_{\Ran}\barB_X(\cA)$, with strict model
$\cA^!=(\cA^i)^\vee$ when finite Koszul formality holds.
\end{remark}

\subsection{Verification of coalgebra axioms}

\begin{theorem}[Bar coalgebra and post-Verdier algebra; \ClaimStatusProvedHere]\label{thm:coalgebra-via-NAP}
Under the hypotheses of Construction~\ref{const:A-dual-intrinsic},
the direct-sum bar object $\barB_X^{\mathrm{ch}}(\mathcal{A})$ is a
well-defined conilpotent chiral coalgebra. In the completed
bar-length ambient it is complete pro-conilpotent. In the strict
direct-sum case:
\begin{enumerate}
\item \emph{Coassociativity:}
$(\Delta \otimes \operatorname{id}) \circ \Delta = (\operatorname{id} \otimes \Delta) \circ \Delta$.
\item \emph{Counit:}
$(\epsilon \otimes \operatorname{id}) \circ \Delta = \operatorname{id} = (\operatorname{id} \otimes \epsilon) \circ \Delta$.
\item \emph{Coderivation (graded):}
$\Delta(d(c)) = \sum d(c_{(1)}) \otimes c_{(2)} + (-1)^{|c_{(1)}|} c_{(1)} \otimes d(c_{(2)})$.
\item \emph{Conilpotency:}
For each homogeneous bar element $b$, there exists $N$ such that
$\Delta^{(N)}(b) = 0$.
\end{enumerate}
Its Verdier dual, respectively continuous Verdier dual,
$\mathcal{A}^!_\infty=\mathbb{D}_{\Ran}\barB_X^{\mathrm{ch}}(\mathcal{A})$
is a factorization algebra whose multiplication, unit, and
differential are the Verdier duals of these coalgebra operations.
\end{theorem}

\begin{proof}[Proof via geometric stratifications]
The coproduct
$\Delta\colon\barB_X^{\mathrm{ch}}(\mathcal{A})\to
\barB_X^{\mathrm{ch}}(\mathcal{A})\otimes
\barB_X^{\mathrm{ch}}(\mathcal{A})$
arises from the geometric correspondence:
\[C_k(X) \to \bigcup_{I \sqcup J = [k]} C_{|I|}(X) \times C_{|J|}(X)\]
that splits points into two groups.

The composition $(\Delta \otimes \text{id}) \circ \Delta$ corresponds to:
\[C_k(X) \to C_{|I|}(X) \times C_{|J|}(X) \times C_{|K|}(X)\]
for $I \sqcup J \sqcup K = [k]$.

The two bracketings give the same ordered triple decomposition
$I \sqcup J \sqcup K=[k]$, giving coassociativity.

The counit
$\epsilon\colon\barB_X^{\mathrm{ch}}(\mathcal{A})\to\omega_X$
corresponds to the projection:
\[C_k(X) \to X\]
selecting a single point (say, the first).

The composition $(\epsilon \otimes \text{id}) \circ \Delta$ corresponds to:
\[C_k(X) \xrightarrow{\Delta} C_1(X) \times C_{k-1}(X) \xrightarrow{\epsilon \times \text{id}} X \times C_{k-1}(X) \cong C_{k-1}(X)\]

This is the identity on $C_{k-1}(X)$, verifying the counit axiom.

The bar differential is the sum of internal differential, de~Rham
differential, and residue along collision divisors. Geometrically,
the residue term corresponds to:
\[d: C_k(X) \to \bigcup_{i<j} D_{ij} \times C_{k-2}(X)\]
where $D_{ij}$ is the collision divisor.

The coderivation property:
\[\Delta \circ d = (d \otimes \text{id} + \text{id} \otimes d) \circ \Delta\]

follows from the identity:
\[\text{(split then collide)} = \text{(collide then split on left)} + \text{(collide then split on right)}\]

This is the combinatorial identity for configuration space strata, which holds because boundary divisors satisfy:
\[\partial(\overline{C}_k(X)) = \bigcup_{\text{strata}} D_{\sigma}\]
with compatible orientations.

Conilpotency is the bar-length filtration in the direct-sum
category. An element represented on $C_k(X)$ has at most $k$
collision factors, so every iterated reduced coproduct of length
$>k$ is zero. In the completed category this becomes
pro-conilpotence. Applying $\mathbb{D}_{\Ran}$, or the continuous
dual in the completed case, converts these coalgebra identities into
the factorization-algebra identities for $\mathcal{A}^!_\infty$.
\end{proof}

\section{The bar coalgebra, its cohomology, and the Verdier-dual algebra}

\begin{theorem}[Bar coalgebra, Koszul-dual coalgebra, and Verdier-dual algebra; \ClaimStatusProvedHere]\label{thm:bar-computes-dual}
\index{Verdier duality!bar coalgebra and dual algebra|textbf}
\index{non-abelian Poincare duality}
\textup{[Regime: quadratic finite-type, with completed continuous
variant
\textup{(}Convention~\textup{\ref{conv:regime-tags})}.]}

Let $\mathcal{A}$ be an augmented chiral algebra on a smooth curve $X$,
valued in holonomic $\mathcal{D}$-modules
(Convention~\ref{subsec:ambient-category}). Assume that the configuration-space
pushforwards $\pi_*(\mathcal{A}^{\boxtimes(n+1)} \otimes \omega^{\log}_{C_{n+1}(X)})$
are well-defined in the derived category of holonomic
$\mathcal{D}$-modules on $X$, that the relevant de~Rham cohomology
groups are finite-dimensional after imposing compact support on the
open strata, and that the orientation line of the curve has been
fixed. In the non-finite case replace the bar object by its separated
complete pro-conilpotent completion and replace all duals by
continuous Verdier duals.

Then:
\begin{enumerate}[label=\textup{(\alph*)}]
\item $\bar{B}^{\mathrm{ch}}(\mathcal{A})$ is a conilpotent
 factorization coalgebra.
\item Its cohomology
 $\mathcal{A}^i:=H^\bullet\bar{B}^{\mathrm{ch}}(\mathcal{A})$
 is the Koszul-dual coalgebra on the finite Koszul locus.
\item Its Verdier dual
 $\mathcal{A}^!_\infty:=
 \mathbb{D}_{\Ran}\bar{B}^{\mathrm{ch}}(\mathcal{A})$
 is a factorization algebra, equipped with the perfect finite-type
 evaluation pairing, respectively the continuous evaluation pairing,
 \[
 \bar{B}^{\mathrm{ch}}(\mathcal{A})\otimes
 \mathcal{A}^!_\infty \longrightarrow \omega_{\Ran}.
 \]
\item If $\mathcal{A}$ is finite Koszul and the transferred
 post-Verdier $A_\infty$-structure is formal, then
 \[
 \mathcal{A}^!_\infty \simeq \mathcal{A}^!
 :=(\mathcal{A}^i)^\vee
 \]
 as factorization algebras.
\end{enumerate}
No morphism
$\bar{B}^{\mathrm{ch}}(\mathcal{A})\to\mathcal{A}^!$
is asserted. The former is a coalgebra, the latter an algebra.
\end{theorem}

\begin{proof}
The holonomicity hypothesis puts every bar component in the
bounded derived category of holonomic $\mathcal{D}$-modules on
$\overline{C}_{n+1}(X)$. Definition~\ref{def:geometric-bar}
therefore gives a conilpotent factorization coalgebra, with
coproduct induced by splitting configurations.

Apply $\mathbb{D}_{\Ran}$ on the Verdier-dualizable holonomic
subcategory. Verdier duality is involutive there and exchanges
$j_*$ with $j_!$
(Lemma~\ref{lem:verdier-extension-exchange}). Hence the dual of a
bar component is represented on the open configuration space by
distributional sections, and the canonical evaluation pairing is
perfect in the finite-type case and continuous after completion.
The differential compatibility is componentwise:
\begin{align*}
\mathbb{D}(d_{\text{int}}) &= d_{\text{int}}^{\text{dual}} \quad \text{(internal differential)}\\
\mathbb{D}(d_{\text{res}}) &= d_{\delta} \quad \text{(delta function insertion)}\\
\mathbb{D}(d_{\text{dR}}) &= d_{\text{dR}}^{\text{dist}} \quad \text{(de Rham on distributions)}
\end{align*}
so the dual object is the factorization algebra
$\mathcal{A}^!_\infty$ of Construction~\ref{const:A-dual-intrinsic}.
The coalgebra coproduct on $\barB^{\mathrm{ch}}(\mathcal{A})$
duals to the multiplication on $\mathcal{A}^!_\infty$, and the
counit duals to the unit.

In the finite-type case Verdier duality gives, with the dimension
shifts fixed by Theorem~\ref{thm:verdier-config},
\[
H^\bullet(\mathcal{A}^!_\infty)
\cong
H^\bullet(\bar{B}^{\mathrm{ch}}(\mathcal{A}))^\vee.
\]
In the completed case the right hand side is the continuous dual and
requires the Mittag--Leffler condition for cohomology to commute with
the inverse limit. On the finite Koszul locus the right hand side is
$(\mathcal{A}^i)^\vee=\mathcal{A}^!$. Formality of the transferred
$A_\infty$-structure, when available, identifies the homotopy dual
algebra $\mathcal{A}^!_\infty$ with its strict model $\mathcal{A}^!$.
\end{proof}

\begin{example}[Heisenberg specialization]\label{ex:bar-dual-heisenberg}
For $\cA = \cH_k$, Theorem~\ref{thm:bar-computes-dual} gives
$H^\bullet\bar{B}^{\mathrm{ch}}(\cH_k)=\cH_k^i$ and
$\mathbb{D}_{\Ran}\bar{B}^{\mathrm{ch}}(\cH_k)
\simeq(\cH_k)^!_\infty$. The bar complex is built from
logarithmic forms on $\overline{C}_n(X)$; its Verdier dual is the
factorization algebra represented by distributional currents on
$C_n(X)$. Since $\cH_k$ is Koszul, the homotopy dual is formal and
$(\cH_k)^!_\infty\simeq\cH_k^!
= \mathrm{Sym}^{\mathrm{ch}}(V^*)$
(cf.\ \S\ref{sec:frame-inversion}).
\end{example}

\subsection{Explicit computation in low degrees}

\begin{computation}[Degree 0 and 1; \ClaimStatusProvedHere]\label{comp:bar-dual-low-degrees}
For a finite-type quadratic chiral algebra
$\mathcal{A}=T_{\mathrm{ch}}(V)/(R)$, the low arities display the
coalgebra, its Verdier dual algebra, and the residue pairing.

\emph{Degree 0.}
\[\bar{B}^0(\mathcal{A}) = \mathcal{A}\]
\[(\mathcal{A}^!_\infty)^{(0)} = \mathbb{D}(\mathcal{A}) \otimes \omega_X\]

There is a canonical evaluation pairing
\[
\mathcal{A}\otimes \mathbb{D}(\mathcal{A})\longrightarrow\omega_X .
\]
A map
$\Phi_0\colon\mathcal{A}\to\mathbb{D}(\mathcal{A})\otimes\omega_X$
exists only after choosing a perfect self-duality or orientation
trivialization of $\mathcal{A}$. The bar/Verdier construction does
not supply such a Calabi--Yau pairing by itself.

\emph{Degree 1.}
\[\bar{B}^1(\mathcal{A}) = \Gamma(\overline{C}_2(X), \mathcal{A}^{\boxtimes 2} \otimes \eta_{12})\]

The dual is:
\[(\mathcal{A}^!_\infty)^{(1)}
= \mathcal{D}^*(C_2(X), (\mathbb{D}_X\mathcal{A})^{\boxtimes 2})\]

The Verdier-dual representative is the boundary-localization class
\[
\Phi_1(\phi_1 \otimes \phi_2 \otimes \eta_{12})
=
\mathbb{D}(\phi_1 \otimes \phi_2)\otimes [D_{12}]_{\mathrm{BM}},
\]
often denoted analytically by
$\mathbb{D}(\phi_1\otimes\phi_2)\otimes\delta(z_1-z_2)$.
Here $\eta_{12}=d\log(z_1-z_2)$ is the logarithmic form and
$[D_{12}]_{\mathrm{BM}}$ is the Verdier-dual Borel--Moore class of
the collision divisor.

The pairing is given by the residue:
\[\langle \eta_{12}, \delta_{12} \rangle = \operatorname{Res}_{z_1 = z_2}\, d\log(z_1 - z_2) = 1\]

(The pointwise product $\frac{1}{z}\delta(z)$ is not defined in distribution theory;
the correct formulation uses the residue pairing between logarithmic forms and
delta currents supported on the diagonal.) This is the fundamental Verdier pairing.

\emph{Degree 2 (first nontrivial case).}
\[\bar{B}^2(\mathcal{A}) = \Gamma(\overline{C}_3(X), \mathcal{A}^{\boxtimes 3} \otimes (\eta_{12} \wedge \eta_{23} + \text{cyc}))\]

The dual is:
\[(\mathcal{A}^!_\infty)^{(2)}
= \mathcal{D}^*(C_3(X),(\mathbb{D}_X\mathcal{A})^{\boxtimes 3})\]

The corresponding representative is the iterated localization class
of the collision flag $D_{12}\cap D_{23}$:
\[
\Phi_2(\phi_1 \otimes \phi_2 \otimes \phi_3 \otimes
\eta_{12} \wedge \eta_{23})
=
\mathbb{D}(\phi_1 \otimes \phi_2 \otimes \phi_3)
\otimes [D_{12}\cap D_{23}]_{\mathrm{BM}}^{\mathrm{flag}}.
\]
The symbol $\delta(z_1-z_2)\delta(z_2-z_3)$ is only a mnemonic for
this flag class.

The Arnold relations on the bar side:
\[\eta_{12} \wedge \eta_{23} + \eta_{23} \wedge \eta_{31} + \eta_{31} \wedge \eta_{12} = 0\]

correspond, under Verdier duality, to relations among constructible
sheaves supported on the diagonal strata $\{z_i=z_j\}$. The precise
formulation requires $\mathcal{D}$-modules, perverse sheaves, or
Borel--Moore chains on the stratified compactification. Literal
products of delta distributions have the wrong support data. The naive products
$\delta(z_1{-}z_2)\delta(z_2{-}z_3)$,
$\delta(z_2{-}z_3)\delta(z_3{-}z_1)$, and
$\delta(z_3{-}z_1)\delta(z_1{-}z_2)$ all collapse to the same full
diagonal support and forget the flag orientation data; they therefore
do not carry the Arnold relation. The correct dual statement is that
the Verdier dual of the Arnold relation in $H^2(C_3(X))$ is the
corresponding Borel--Moore relation in the shifted Verdier-dual
degree.
\end{computation}

\section{Koszul pairs and symmetric duality}

\begin{proposition}[Verdier compatibility in a chiral Koszul pair; \ClaimStatusConditional]\label{prop:koszul-pair-NAP}
\index{Koszul pair!via non-abelian Poincare duality}
Let $(\mathcal{A}_1,\mathcal{A}_2)$ be two chiral algebras on $X$
equipped with the antecedent twisting morphisms, compatible
filtrations, conilpotent finite or completed Koszul coalgebras, and
finite-type or continuous dualizability required in
Definition~\ref{def:chiral-koszul-pair}. The Verdier-compatibility
part of the chiral Koszul-pair datum is the pair of identifications
\begin{align}
\mathbb{D}_{\Ran}\bar{B}^{\mathrm{ch}}(\mathcal{A}_1)
&\xrightarrow{\sim} \mathcal{A}_2,\\
\mathbb{D}_{\Ran}\bar{B}^{\mathrm{ch}}(\mathcal{A}_2)
&\xrightarrow{\sim} \mathcal{A}_1.
\end{align}
Equivalently, on the finite Koszul locus,
\[
\mathcal{A}_i^i
=H^\bullet\bar{B}^{\mathrm{ch}}(\mathcal{A}_i),
\qquad
\mathcal{A}_i^!=(\mathcal{A}_i^i)^\vee,
\]
and the notation $\mathcal{A}_2\simeq\mathcal{A}_1^!$,
$\mathcal{A}_1\simeq\mathcal{A}_2^!$ names the post-Verdier
algebras, not the bar coalgebras.

\emph{NAP shadow.}
If the locally constant topological shadows of these chiral
algebras satisfy the Ayala--Francis compactness, orientation,
support, and dualizability hypotheses, factorization homology gives
\[\int_X \mathcal{A}_1 \simeq \mathbb{D}\left(\int_{-X} \mathcal{A}_2\right)\]

with $\int_X$ the factorization homology, $-X$ the orientation
reversal, and $\mathbb{D}$ Verdier duality. This topological shadow
does not replace the chiral twisting and filtration data.
\end{proposition}

\begin{proof}
Theorem~\ref{thm:bar-computes-dual} separates the two assertions:
$\bar{B}^{\mathrm{ch}}(\mathcal{A})$ is the coalgebra and
$\mathbb{D}_{\Ran}\bar{B}^{\mathrm{ch}}(\mathcal{A})$ is the
post-Verdier algebra. Definition~\ref{def:chiral-koszul-pair}
requires twisting morphisms, filtrations, Koszulness, and precisely
the Verdier-compatible identifications of these post-Verdier algebras
with the opposite algebra in the pair. Applying factorization
homology gives the displayed NAP equivalence only after the
Ayala--Francis hypotheses are imposed on the locally constant
topological shadows.
\end{proof}

\begin{theorem}[Symmetric Koszul duality; \ClaimStatusConditional]\label{thm:symmetric-koszul}
\textup{[Regime: quadratic finite-type on the Koszul locus, with
completed continuous variant
\textup{(}Convention~\textup{\ref{conv:regime-tags})}.]}

If $(\mathcal{A}_1, \mathcal{A}_2)$ is a finite-type chiral Koszul
pair with Verdier-dualizable bar components and formal post-Verdier
dual algebras, then:
\begin{align}
H^\bullet\bar{B}^{\mathrm{ch}}(\mathcal{A}_1)
&\simeq \mathcal{A}_1^i,\\
H^\bullet\bar{B}^{\mathrm{ch}}(\mathcal{A}_2)
&\simeq \mathcal{A}_2^i,\\
\mathbb{D}_{\Ran}\bar{B}^{\mathrm{ch}}(\mathcal{A}_1)
&\simeq \mathcal{A}_2\simeq\mathcal{A}_1^!,\\
\mathbb{D}_{\Ran}\bar{B}^{\mathrm{ch}}(\mathcal{A}_2)
&\simeq \mathcal{A}_1\simeq\mathcal{A}_2^!,\\
\Omega^{\mathrm{ch}}\bar{B}^{\mathrm{ch}}(\mathcal{A}_i)
&\simeq \mathcal{A}_i \qquad (i=1,2).
\end{align}

Thus $\Omega^{\mathrm{ch}}\bar{B}^{\mathrm{ch}}$ is inversion on each
member, while $\mathbb{D}_{\Ran}\bar{B}^{\mathrm{ch}}$ is the duality
operation between the two members. In the completed case all four
displayed duality lines are read with continuous duals and completed
cobar.
\end{theorem}

\begin{proof}
The cohomology identifications are the definition of
$\mathcal{A}_i^i$ on the finite Koszul locus. The Verdier
identifications are exactly the compatibility data of
Definition~\ref{def:chiral-koszul-pair}, rewritten using
Construction~\ref{const:A-dual-intrinsic}. The last line is the
bar-cobar counit on the Koszul locus
(Theorem~\ref{thm:bar-cobar-adjunction} and
Theorem~\ref{thm:bar-cobar-isomorphism-main}). It reconstructs the
same algebra from its own bar coalgebra.
\end{proof}

\begin{example}[Heisenberg specialization]\label{ex:symmetric-heisenberg}
For the Koszul pair $(\cH_k, \mathrm{Sym}^{\mathrm{ch}}(V^*))$,
Theorem~\ref{thm:symmetric-koszul} gives:
$\mathbb{D}_{\Ran}\bar{B}^{\mathrm{ch}}(\cH_k)
\simeq \mathrm{Sym}^{\mathrm{ch}}(V^*)$
and
$\Omega^{\mathrm{ch}}\bar{B}^{\mathrm{ch}}(\cH_k)\simeq\cH_k$.
This specializes to the Lie--Com exchange: the chiral Lie structure
of $\cH_k$ dualizes to the chiral commutative
structure of $\mathrm{Sym}^{\mathrm{ch}}(V^*)$, and vice versa
(\S\ref{sec:frame-koszul-dual}).
\end{example}

\begin{remark}[Surface observables via factorization homology]
\label{rem:surface-observables-fh}
\index{factorization homology!surface observables}
\index{Ayala--Francis!Poincare--Koszul duality@Poincar\'e--Koszul duality}
For a closed genus-$g$ surface $\Sigma_g$, the factorization
homology
$H_{\Sigma_g}(\cA) := \int_{\Sigma_g} \cA$
models the topological factorization-homology observables of the
locally constant shadow of~$\cA$. It is not, without an additional
comparison theorem, the chiral Hochschild complex
$\ChirHoch^\bullet(\cA,\cA)$ or the categorical Hochschild complex of
a module category. Under the Ayala--Francis compactness,
orientation, support, and dualizability hypotheses,
Poincar\'e/Koszul duality~\cite{AF15} gives
\[
\int_{\Sigma_g} \cA
\;\simeq\;
\mathbb{D}\!\int_{-\Sigma_g} \cA^!,
\]
which is the topological shadow of the modular bar-cobar story:
the left side is factorization homology with~$\cA$; the right
side is factorization cohomology with the post-Verdier dual
algebra~$\cA^!$, not with the bar coalgebra. At genus~$0$ this is
the topological shadow of Theorem~A. In the modular chiral lane at
$g \geq 1$, the corresponding bar-family observable is
$R\Gamma(\overline{\mathcal{M}}_g,\barB_g(\cA))$ after the
proper-support and completion hypotheses are imposed. The modular
characteristic package extracts numerical invariants from that
derived object.
\end{remark}

\section{Non-quadratic cases and completion}

\begin{remark}[Necessity of completion]\label{rem:why-completion}
For non-quadratic chiral algebras (W-algebras, Yangians, etc.), the bar construction produces infinitely many generators in each degree:

The bar complex is not finitely generated. The cohomological
coalgebra $\mathcal{A}^i=H^\bullet\barB(\mathcal{A})$ therefore
lives naturally in a completed conilpotent category, and the algebra
$\mathcal{A}^!=(\mathcal{A}^i)^\vee$ is a continuous Verdier-linear
dual.

The Beilinson--Drinfeld remedy is $I$-adic completion:
\[\widehat{\bar{B}}(\mathcal{A}) = \varprojlim_n \bar{B}(\mathcal{A})/I^n\]
where $I = \ker(\epsilon: \bar{B}(\mathcal{A}) \to \mathbb{C})$ is the coaugmentation ideal.

\emph{Geometric interpretation.} The completion is the inverse limit
\[\widehat{\bar{B}}(\mathcal{A}) = \varprojlim_n \bar{B}(\mathcal{A})/F^n\]
over the conilpotent filtration by bar degree (see Appendix~\ref{app:nilpotent-completion}), which organizes the bar complex by the number of collision points on~$\overline{C}_n(X)$.
\end{remark}

\begin{definition}[Completed dual coalgebra and dual algebra]\label{def:completed-dual}
For a chiral algebra $\mathcal{A}$, define
\[
\widehat{\mathcal{A}^i}
:=
H^\bullet\widehat{\bar{B}}^{\mathrm{ch}}(\mathcal{A}),
\qquad
\widehat{\mathcal{A}^!}
:=
(\widehat{\mathcal{A}^i})^\vee_{\mathrm{cont}}.
\]
The completed homotopy dual algebra is
\[
\widehat{\mathcal{A}^{!}_{\infty}}
:=
\mathbb{D}^{\mathrm{cont}}_{\Ran}
\widehat{\bar{B}}^{\mathrm{ch}}(\mathcal{A}).
\]
The completion is taken with respect to the conilpotent
bar-degree filtration, equivalently the powers of the coaugmentation
ideal.
\end{definition}

\begin{theorem}[Completion and Koszul duality; \ClaimStatusConditional]\label{thm:completion-koszul}
For chiral algebras satisfying:
\begin{enumerate}
\item Finite generation over $\mathcal{D}_X$
\item Polynomial growth of structure constants
\item Separated and complete bar-degree filtration
\item Finite-dimensional graded quotients with Mittag--Leffler
 transition maps on cohomology
\item Continuous Verdier dualizability with proper support on every
 Fulton--MacPherson stratum
\end{enumerate}

The completed bar construction is a complete conilpotent chiral
coalgebra whose cohomology is the completed dual coalgebra:
\[
H^\bullet\widehat{\bar{B}}^{\mathrm{ch}}(\mathcal{A})
\cong \widehat{\mathcal{A}^i}.
\]
Its continuous Verdier dual is the completed homotopy dual algebra:
\[
\mathbb{D}^{\mathrm{cont}}_{\Ran}
\widehat{\bar{B}}^{\mathrm{ch}}(\mathcal{A})
\simeq
\widehat{\mathcal{A}^{!}_{\infty}}.
\]

The completed cobar counit:
\[
\Omega^{\mathrm{ch}}\widehat{\bar{B}}^{\mathrm{ch}}(\mathcal{A})
\xrightarrow{\sim} \widehat{\mathcal{A}}
\]
recovers the completed algebra on the completed Koszul surface. It
recovers $\mathcal{A}$ itself only when the completion map
$\mathcal{A}\to\widehat{\mathcal{A}}$ is an equivalence.
\end{theorem}

\begin{proof}
\emph{Step 1: Verdier duality extends to completed sheaves.}
The bar complex carries a filtration
$F_n \bar{B}^{\text{ch}}(\mathcal{A})$ by bar degree $\leq n$.
Since $\mathcal{A}$ is finitely generated over $\mathcal{D}_X$
(hypothesis (1)), each $\bar{B}^{\text{ch}}_n(\mathcal{A})$ is a
holonomic $\mathcal{D}$-module on $\overline{C}_n(X)$: the bar
construction at degree~$n$ involves $n$-fold chiral tensor products
of~$\mathcal{A}$, which are holonomic by closure of holonomic
$\mathcal{D}$-modules under tensor product and proper pushforward.
For a complete inverse system
$\mathcal{F}=\varprojlim_n\mathcal{F}/F^n$, continuous Verdier
duality is defined by
\[
\mathbb{D}^{\mathrm{cont}}(\mathcal{F})
:=\varinjlim_n\mathbb{D}(\mathcal{F}/F^n).
\]
This is a completed duality convention controlled by the topology;
it is not an unconditional commutation rule for arbitrary inverse
limits.

\emph{Step 2: Completion filtration is compatible with Verdier duality.}
The coaugmentation ideal
$I=\ker(\epsilon:\bar{B}^{\text{ch}}\to\mathbb{C})$ defines the
$I$-adic filtration. Each quotient $\bar{B}^{\text{ch}}/I^n$
involves only configuration spaces $C_k(X)$ with $k\leq n$, hence is
geometric. Polynomial growth (hypothesis (2)) controls the size of
the graded pieces; it does not by itself imply Mittag--Leffler.
The Mittag--Leffler condition is hypothesis (4), and gives
$\varprojlim^1=0$ for the cohomology inverse system.

\emph{Step 3: Lifting to completions.}
By Theorem~\ref{thm:bar-computes-dual}, each finite quotient
$\bar{B}^{\mathrm{ch}}(\mathcal{A})/I^n$ is a conilpotent
coalgebra and its Verdier dual is an algebra. The universal
property of the inverse limit, together with the Mittag-Leffler
condition from Step~2, yields:
\[
\widehat{\bar{B}}^{\mathrm{ch}}(\mathcal{A})
=\varprojlim_n \bar{B}^{\mathrm{ch}}(\mathcal{A})/I^n
\]
as a complete coalgebra, and
\[
\mathbb{D}^{\mathrm{cont}}_{\Ran}
\widehat{\bar{B}}^{\mathrm{ch}}(\mathcal{A})
\simeq
\varinjlim_n
\mathbb{D}_{\Ran}\!\left(
\bar{B}^{\mathrm{ch}}(\mathcal{A})/I^n
\right)
\]
as its continuous Verdier-dual algebra. Taking cohomology gives
$H^\bullet\widehat{\bar{B}}^{\mathrm{ch}}(\mathcal{A})
\cong\widehat{\mathcal{A}^i}$, and continuous duality gives
$H^\bullet\widehat{\mathcal{A}^{!}_{\infty}}
\cong\widehat{\mathcal{A}^!}$.

\emph{Step~4: Cobar.}
By hypothesis~(1), each bar component is holonomic
(Step~1), so $\dim \mathrm{CH}^n(\mathcal{A}) < \infty$
(de~Rham cohomology of holonomic $\mathcal{D}$-modules is
finite-dimensional). Hence the cobar construction
$\Omega^{\mathrm{ch}}\widehat{\bar{B}}^{\mathrm{ch}}(\mathcal{A})$
is well-defined and carries a complete filtration compatible with
the differential. The bar-cobar adjunction counit
$\Omega(\bar{B}(\mathcal{A})) \to \mathcal{A}$
(Theorem~\ref{thm:bar-cobar-adjunction}) is a quasi-isomorphism
on the Koszul locus
(Theorem~\ref{thm:bar-cobar-isomorphism-main}).
In the completed setting, all operations are continuous with
respect to the $I$-adic topology and the Mittag-Leffler
condition (Step~2) ensures $\varprojlim^1 = 0$, so the
completed counit
$\Omega^{\mathrm{ch}}\widehat{\bar{B}}^{\mathrm{ch}}(\cA)
\to \widehat{\cA}$ is a quasi-isomorphism. If
$\cA\simeq\widehat{\cA}$, this is the stated reconstruction of
$\cA$.
\end{proof}

\subsection{Application: $\mathcal{W}$-algebras}

\begin{example}[First completed window for \texorpdfstring{$W_3$}{W3}]\label{ex:W-completion}
For the $W_3$ algebra with generators $T(z) = L(z)$ (weight 2) and $W(z)$ (weight 3):

\emph{The composite field.}
\[\Lambda = {:}TT{:} - \frac{3}{10}\,\partial^2 T\]
is the quasi-primary composite of weight 4. It is absent from the
strong generating set of~$W_3$ and appears in
$\bar{B}^{\mathrm{ch}}(W_3)$ with OPE coefficient
$\frac{16}{22+5c}$ in the $W$--$W$ channel
(Definition~\ref{def:lambda-complete}).

\emph{Completed dual coalgebra, first associated-graded window.}
\[\widehat{W_3^i} = \text{Free coalgebra}(L^*, W^*, \Lambda^*, (\Lambda^*)^{(2)}, (\Lambda^*)^{(3)}, \ldots)\]

where $(\Lambda^*)^{(n)}$ are descendant towers. This display records
the completed associated graded generated by the first composite
field and its descendants; the full completed coalgebra also carries
the topology and higher residue terms determined by the complete
$W$--$W$ OPE.

\emph{Coproduct.}
\begin{align*}
\Delta(L^*) &= 0 \quad \text{(primitive)}\\
\Delta(W^*) &= 0 \quad \text{(primitive)}\\
\Delta(\Lambda^*) &= L^* \otimes L^* - \frac{3}{10}\left(\partial^2 L^* \otimes 1 + 1 \otimes \partial^2 L^*\right) \quad \text{(composite)}
\end{align*}

The completed dual algebra is the continuous Verdier-linear dual
\[
\widehat{W_3^!}:=(\widehat{W_3^i})^\vee_{\mathrm{cont}}.
\]

\emph{The differential.}
The differential on $\Lambda^*$ encodes the $W$--$W$ OPE:
since $W(z)W(w) \sim \frac{16}{22+5c}\frac{\Lambda(w)}{(z-w)^2} + \cdots$, the bar differential acts by
\[d(\Lambda^*) = \frac{16}{22+5c}\, W^* \otimes W^*\]
with the remaining terms determined by the full $W$--$W$ OPE (see Definition~\ref{def:lambda-complete} and Computation~\ref{comp:w3-bar-degree2} for the explicit structure constants).

\emph{Cobar reconstruction.}
\[\Omega^{\text{ch}}(\widehat{W_3^i}) \simeq \widehat{W_3}\]

on the completed Koszul surface. The relation
$\Lambda = :T \cdot T: - \frac{3}{10}\partial^2 T$ (with OPE
coefficient $\frac{16}{22+5c}$) contributes the displayed
Maurer--Cartan component in the cobar complex.
\end{example}

\section{Verdier resolution of the type obstruction}

\begin{theorem}[Verdier resolution of the type obstruction; \ClaimStatusConditional]\label{thm:main-NAP-resolution}
Assume the finite-type hypotheses of
Theorem~\ref{thm:bar-computes-dual}, or the completed hypotheses of
Theorem~\ref{thm:completion-koszul}.
Construction~\ref{const:A-dual-intrinsic} defines the post-Verdier
homotopy dual algebra
$\cA^!_\infty=\mathbb{D}_{\Ran}\barB_X^{\mathrm{ch}}(\cA)$ with no
pre-existing dual partner or orthogonal-relation presentation.
Theorem~\ref{thm:bar-computes-dual} identifies
$\cA^i=H^\bullet\barB_X^{\mathrm{ch}}(\cA)$ as the Koszul-dual
coalgebra on the finite Koszul locus and
$\cA^!=(\cA^i)^\vee$ as the strict Verdier-linear dual algebra.
For a Koszul pair $(\mathcal{A}_1,\mathcal{A}_2)$,
Theorem~\ref{thm:symmetric-koszul} gives
\[
\mathbb{D}_{\Ran}\bar{B}^{\mathrm{ch}}(\mathcal{A}_1)
\simeq \mathcal{A}_2\simeq\mathcal{A}_1^!,
\qquad
\mathbb{D}_{\Ran}\bar{B}^{\mathrm{ch}}(\mathcal{A}_2)
\simeq \mathcal{A}_1\simeq\mathcal{A}_2^!,
\]
and the cobar equivalences
$\Omega^{\mathrm{ch}}\barB^{\mathrm{ch}}(\mathcal{A}_i)
\simeq\mathcal{A}_i$ are bar-cobar inversion.
\end{theorem}

\begin{proof}
The construction of $\cA^!_\infty$ is
Construction~\ref{const:A-dual-intrinsic}. The coalgebra and
algebra identifications are Theorem~\ref{thm:bar-computes-dual}.
The symmetric pair statement is Theorem~\ref{thm:symmetric-koszul},
which is the object-level form of the NAP equivalence
$\int_X \mathcal{A}_1
\simeq \mathbb{D}\bigl(\int_{-X}\mathcal{A}_2\bigr)$
(Theorem~\ref{thm:verdier-NAP}). The final clause is the
bar-cobar counit on the Koszul or completed Koszul surface.
\end{proof}

The completed bar construction (Theorem~\ref{thm:completion-koszul})
extends this to non-quadratic algebras only after the separated
completion, Mittag--Leffler, proper-support, and continuous-duality
hypotheses are imposed. The quantum corrections (curvature from
periods on $\overline{\mathcal{M}}_g$) appear in
\S\ref{sec:koszul-defects} and Chapter~\ref{chap:higher-genus}.

\begin{remark}[Use in the main theorems]\label{rem:nap-usage}
Theorem~\ref{thm:main-NAP-resolution} supplies the post-Verdier
dual algebra for Theorem~A$_2$
(Theorem~\ref{thm:bar-cobar-isomorphism-main});
Theorem~\ref{thm:verdier-bar-cobar} supplies the Verdier
compatibility input for Theorem~C
(Theorem~\ref{thm:quantum-complementarity-main}); and the
explicit NAP computations of
Chapter~\ref{chap:NAP-computations} independently verify
Heisenberg, Kac--Moody, and $\beta\gamma$/bc Koszul pairs.
\end{remark}

\section{K3 Poincar\'e-duality shadows}
\label{sec:pd-supplement}

\begin{remark}[Mukai pairing as Poincar\'e duality]
\label{rem:pd-1}
The Mukai pairing on
$\Lambda_{\mathrm{Muk}}=\mathrm{II}_{4,20}$ is the
Poincar\'e pairing on K3 cohomology transported to the Mukai lattice:
signature $(4,20)$, with $c_+=4$ and $c_-=20$. The scalar
Mukai-doubling value $K^{\kappa_{\mathrm{ch}}}=2c_+=8$ is the
chiral shadow used in the K3 bialgebra package, not a statement that
the chiral algebra $\mathbf H_{\Delta_5}$ is itself the Mukai
lattice. It is also not the Borcherds weight
$\kappa_{\mathrm{BKM}}(\Delta_5)=5$ and not the K3$\times E$
Heisenberg shadow $\kappa_{\mathrm{ch}}^{\mathrm{Heis}}=3$. The
Poincar\'e-duality dimension $24=c_++c_-$ equals the rank of the
Mukai lattice and matches the $24$ Kodaira $I_1$ fibres on
$E^{\mathrm{nod,sm}}_{24}$ in the bi-based Ran pair. The lattice
duality is implemented by the $\mathrm{SO}(4,20;\mathbb Z)$-action
on K3 moduli \cite{Mukai1987,Borcherds1998}; the chiral lift is the
conditional Vol.~III comparison of
Chapter~\ref{V3-ch:k3-chiral-bialgebra-platonic}.
\end{remark}

\begin{remark}[Three faces of Poincar\'e duality at K3]
\label{rem:pd-2}
Poincar\'e duality at K3 has three faces:
(i)~the Mukai or categorical Calabi--Yau pairing on
$D^b\mathrm{Coh}(K3)$ (CY side);
(ii)~the chiral Hochschild pairing
$\ChirHoch^\bullet(\mathbf H_{\Delta_5},\mathbf H_{\Delta_5})$
(chiral side);
(iii)~Siegel-modular $\Phi_{10}$-denominator cusp-form pairing on
$\bar{\mathcal A}_2$ (modular-automorphic side).
These are distinct Hochschild and pairing theories. The categorical
Hochschild theory of $D^b\mathrm{Coh}(K3)$, the chiral Hochschild
complex of $\mathbf H_{\Delta_5}$, and the automorphic pairing agree
only after the conditional comparison functors and normalisations of
Vol.~III are imposed. Under that comparison the scalar shadow is the
specialisation $\hbar^2\cdot K^{\kappa_{\mathrm{ch}}}=-1$, hence
$\hbar^2=-1/8$ when $K^{\kappa_{\mathrm{ch}}}=8$
\cite{Lusztig1990,PasolZagier2013}. This scalar shadow uses the
Mukai conductor; the automorphic weight remains
$\kappa_{\mathrm{BKM}}(\Delta_5)=5$.
\end{remark}

\begin{remark}[Quantum-determinant duality]
\label{rem:pd-3-supplement}
The chiral-Hochschild face~(ii) of
Remark~\ref{rem:pd-2} refines on the
RTT side as follows. On the
RTT-style $T$-matrix $T(u,Z)\in\mathrm{End}(V_{24})\otimes
\mathbf H_{\Delta_5}[[u^{-1}]]$ built from the rational-times-theta
factorisation $R_{\mathrm{Sieg,dyn}}(u,Z) =
R^{\mathrm{rat}}_{\mathrm{Yang}}(u)\cdot\theta^{\mathrm{K3}}(u,Z)$,
the conditional Vol.~III comparison asserts that the Molev-style
quantum determinant on the rank-$24$
Mukai-Heisenberg sub-$T$-matrix is central and equal to
$\mathrm{qdet}\,T(u,Z) = C(u)\cdot\Delta_5(Z)\cdot\mathrm{Id}$, with
$C(u)\in 1+u^{-1}\mathbb C[[u^{-1}]]$ a $Z$-independent normalisation
\cite{Molev2007,GritsenkoNikulin1998,Gritsenko1999}. Face~(iii)'s
Siegel-modular $\Delta_5$-pairing is then the Poincar\'e-dual
shadow, at the quantum determinant level, of face~(ii)'s chiral
Hochschild pairing: the two faces coincide only after applying the
Vol.~III comparison and the rational Yang prefactor $C(u)$
(Chapter~\ref{V3-ch:k3-chiral-bialgebra-platonic},
Theorem~\ref{thm:w15-quantum-determinant}).
\end{remark}

\begin{remark}[BKM vertical vs.~Mukai horizontal duality]
\label{rem:pd-4-supplement}
The Mukai pairing rank-$24$ of Remark~\ref{rem:pd-1} is the
\emph{horizontal} duality datum (lattice $\mathrm{II}_{4,20}$,
Mukai-doubling $K^{\kappa_{\mathrm{ch}}}=8$); the
BKM-Chevalley real-simple-root count of $3$ in
Theorem~\ref{thm:w15-chevalley-bkm-generators} gives the
\emph{vertical} duality datum on the rank-three Lorentzian real-root
core $\Lambda^{2,1}_{\Delta_5}$
\cite{GritsenkoNikulin1997,GritsenkoNikulin1998}. This rank-three
core is not an even-unimodular $\mathrm{II}_{2,1}$ lattice and not
the full K3 Mukai lattice. The apparent rank discrepancy $24\neq3$ is the
Mukai-lattice rank (horizontal) vs.~the BKM real-root rank
(vertical). The averaging morphism of
Vol.~III Chapter~\ref{V3-ch:k3-chiral-bialgebra-platonic}'s
Definition~\ref{V3-def:k3-chiral-bialgebra} intertwines the two via the
conditional composite
$\mathrm{II}_{3,19}\hookrightarrow\mathrm{II}_{4,20}\to
\mathrm{Mukai}(K3)$, realising Poincar\'e duality as a commutative
triangle linking the three faces after the categorical, chiral
Hochschild, and Siegel-modular comparison functors have been fixed.
\end{remark}

\begin{remark}[$N$-parametric refinement of the three-faces triangle]
\label{rem:pd-5-N-parametric}
The $N$-parametric family
$\{\mathbf H^{(N)}_{\Delta_5}\}_{N\ge 2}$ of class-$\mathcal S$
K3 chiral bialgebras
(Vol.~I Proposition~\ref{prop:thqg-occ-CD-ANm1-24}) conjecturally
refines the three-face triangle of Remark~\ref{rem:pd-2} to an
$N$-indexed family of commutative triangles. On each face:
\begin{enumerate}
\item \emph{Mukai face.} At $N\ge 2$, the Maulik--Okounkov
 stable-envelope Mukai lattice on $\mathrm{Hilb}^\bullet(K3)$
 is expected to acquire an $\mathfrak{sl}_N$-refinement through the
 $\mathrm{SU}(N)^{24}$ flavour action on the $24$ punctures. The
 finite numerical input is the class-$\mathcal S$ pair
 $(n_v(N),n_h(N))$ of
 Vol.~I Proposition~\ref{prop:thqg-occ-CD-ANm1-24}; interpreting
 those counts as an $N$-refined Mukai-doubling rank is an additional
 comparison datum, not a consequence of the K3 Mukai lattice alone.
\item \emph{Chiral-Hochschild face.} The $N$-parametric chiral Hochschild
 complex $C^\star_{\mathrm{ch}}(\mathbf H^{(N)}_{\Delta_5},
 \mathbf H^{(N)}_{\Delta_5})$ carries an
 $\mathfrak{sl}_N$-refined Mukai-doubling; at $N = 2$ it restores
 $K^{\kappa_{\mathrm{ch}}} = 8$; at $N \ge 3$ the Mukai-doubling
 factorises through the determinant
 $\det_N\colon\mathrm{End}(V_N)\to\mathbb C$.
\item \emph{Siegel-modular face.} The $N=2$ primitive denominator is
 $\Delta_5$ of weight $5$. In the $W$-algebraic class-$\mathcal S$
 surface, Theorem~\ref{thm:walgdeep-gaiotto-siegel-weight} gives
 the higher-rank Borcherds weight
 $k_N^{\mathrm{int}}=N+3$, equivalently spin-cover weight
 $k_N^{\mathrm{spin}}=(N+3)/2$. The comparison of that lift with
 the K3 chiral-bialgebra triangle for $N\ge3$ remains conditional
 on the source-recognition functors and automorphic normalisation.
\end{enumerate}
The conditional $N$-parametric Lusztig specialisation
$\hbar_N^2\cdot K^{\kappa_{\mathrm{ch}}}_N = -1$ predicts
$\hbar_N^2 = -1/K^{\kappa}_N$ where
$K^{\kappa}_N$ is the $N$-parametric Mukai-doubling constant
(not an independent parameter from the $\hbar$ of
Remark~\ref{rem:pd-2}: the bridge identity
$\hbar_{N=2} = \hbar$ recovers $\hbar^2 = -1/K^{\kappa}_2 =
-1/8$ at $N = 2$, and for $N\ge 3$ the convention
$\hbar_N$ refers to the same deformation parameter specialised
against the rank-$N$ BKM core, not a new $\hbar$). The open
obligation is the construction of the comparison functor carrying
the class-$\mathcal S$ Borcherds lift to the chiral Hochschild
pairing of $\mathbf H^{(N)}_{\Delta_5}$; without it the equality is
only a scalar prediction.
\end{remark}
